% TODO checklist - Files inspected for this section:
% - README.md (lines 1-2050): Architecture overview, security model, observability stack
% - Q&A.md (lines 1-7907): Q16-22 (security/multi-tenancy), Q23-28 (LLM resilience), Q42-46 (observability/testing)
% - src/security/: Authentication, authorization, JWT handling
% - src/middleware/: Rate limiting, tenant resolution
% - src/observability/: Metrics, tracing, logging setup
% - src/health/: Health check components
% - SECURITY.md: Security policies and practices

\section{Non-Functional Requirements}
\label{sec:nonfunctional-requirements}

This section specifies the non-functional requirements (NFRs) that govern the quality attributes of the Cineca Agentic Platform. These requirements address cross-cutting concerns including security, reliability, performance, and operational characteristics that are essential for enterprise deployment.

%------------------------------------------------------------------------------
\subsection{Security Requirements}
\label{subsec:nfr-security}

Security is paramount for an enterprise AI agent platform that processes sensitive research data and executes privileged operations. The following requirements establish the security posture:

\begin{table}[htbp]
\centering
\caption{Security Non-Functional Requirements}
\label{tab:nfr-security}
\small
\begin{tabularx}{\textwidth}{|l|X|l|}
\hline
\textbf{ID} & \textbf{Requirement Description} & \textbf{Priority} \\
\hline
NFR-SE-01 & The system shall authenticate all API requests using OAuth 2.0 / OpenID Connect (OIDC) with JWT bearer tokens. & Must \\
\hline
NFR-SE-02 & The system shall validate JWT tokens against OIDC provider public keys (JWKS) with caching and automatic refresh. & Must \\
\hline
NFR-SE-03 & The system shall enforce Role-Based Access Control (RBAC) using scope-based permissions embedded in JWT claims. & Must \\
\hline
NFR-SE-04 & The system shall support configurable RBAC scopes including: \texttt{user:me}, \texttt{admin:all}, \texttt{tools:invoke:basic}, \texttt{tools:invoke:all}, \texttt{graph:read}, \texttt{graph:write}. & Must \\
\hline
NFR-SE-05 & The system shall automatically detect and scrub Personally Identifiable Information (PII) from logs and outputs. & Must \\
\hline
NFR-SE-06 & The system shall prevent Server-Side Request Forgery (SSRF) attacks by validating and restricting outbound request targets. & Must \\
\hline
NFR-SE-07 & The system shall prevent Cypher injection attacks through parameterized queries and input validation. & Must \\
\hline
NFR-SE-08 & The system shall mask sensitive data (tokens, credentials) in all log outputs, displaying only prefix and suffix characters. & Must \\
\hline
NFR-SE-09 & The system shall enforce HTTPS for all external communications in production environments. & Must \\
\hline
NFR-SE-10 & The system shall implement defense-in-depth with multiple security layers (authentication, authorization, input validation, output filtering). & Must \\
\hline
\end{tabularx}
\end{table}

\subsubsection{Authentication Architecture}

The authentication subsystem implements the following security controls:

\begin{itemize}
    \item \textbf{Token Validation:} Every request to protected endpoints must include a valid JWT in the \texttt{Authorization: Bearer <token>} header. Tokens are validated against the OIDC provider's JSON Web Key Set (JWKS).
    
    \item \textbf{Claim Extraction:} Standard OIDC claims (\texttt{sub}, \texttt{iss}, \texttt{aud}, \texttt{exp}, \texttt{iat}) are extracted and validated. Custom claims (scopes, tenant\_id) are used for authorization decisions.
    
    \item \textbf{Key Rotation:} The system caches JWKS keys with automatic refresh when key IDs are not found, supporting transparent key rotation by the identity provider.
    
    \item \textbf{Clock Skew Tolerance:} A configurable clock skew tolerance (default: 30 seconds) accommodates minor time synchronization differences between servers.
\end{itemize}

\subsubsection{Authorization Model}

The platform implements a scope-based RBAC model with the following hierarchy:

\begin{table}[htbp]
\centering
\caption{RBAC Scope Hierarchy}
\label{tab:rbac-scopes}
\small
\begin{tabular}{|l|l|p{7cm}|}
\hline
\textbf{Scope} & \textbf{Level} & \textbf{Permissions Granted} \\
\hline
\texttt{user:me} & User & Read/write own resources (runs, sessions, jobs) \\
\hline
\texttt{tools:invoke:basic} & User & Invoke read-only, non-administrative tools \\
\hline
\texttt{tools:invoke:all} & Power User & Invoke all tools including write operations \\
\hline
\texttt{graph:read} & User & Execute read-only Cypher queries \\
\hline
\texttt{graph:write} & Power User & Execute write Cypher queries (CREATE, MERGE, DELETE) \\
\hline
\texttt{admin:all} & Admin & Full administrative access (tenant management, model configuration, system operations) \\
\hline
\end{tabular}
\end{table}

%------------------------------------------------------------------------------
\subsection{Multi-Tenancy Requirements}
\label{subsec:nfr-multitenancy}

The platform must support multiple organizational tenants with strict data isolation:

\begin{table}[htbp]
\centering
\caption{Multi-Tenancy Non-Functional Requirements}
\label{tab:nfr-multitenancy}
\small
\begin{tabularx}{\textwidth}{|l|X|l|}
\hline
\textbf{ID} & \textbf{Requirement Description} & \textbf{Priority} \\
\hline
NFR-MT-01 & The system shall support logical tenant isolation within shared infrastructure (shared-database, shared-schema pattern). & Must \\
\hline
NFR-MT-02 & The system shall resolve tenant context from JWT claims or \texttt{X-Tenant-ID} header on every request. & Must \\
\hline
NFR-MT-03 & All database queries shall be automatically scoped by \texttt{tenant\_id} to prevent cross-tenant data access. & Must \\
\hline
NFR-MT-04 & The system shall support per-tenant configuration including default models, rate limits, and feature flags. & Must \\
\hline
NFR-MT-05 & The system shall enforce per-tenant quotas for API requests, token consumption, and storage usage. & Should \\
\hline
NFR-MT-06 & The system shall provide tenant-scoped metrics and audit logs for compliance reporting. & Must \\
\hline
NFR-MT-07 & The system shall prevent tenant ID spoofing through cryptographic validation of tenant claims. & Must \\
\hline
NFR-MT-08 & The system shall support tenant onboarding and offboarding through administrative APIs. & Must \\
\hline
\end{tabularx}
\end{table}

\subsubsection{Tenant Isolation Architecture}

The multi-tenancy architecture implements the following isolation mechanisms:

\begin{enumerate}
    \item \textbf{Request-Level Isolation:} Tenant context is established at the middleware layer and propagated through the entire request lifecycle via context variables.
    
    \item \textbf{Query-Level Isolation:} All repository methods automatically append \texttt{WHERE tenant\_id = :tid} predicates. For Cypher queries, tenant isolation is injected during query rewriting.
    
    \item \textbf{Resource-Level Isolation:} All domain entities (Agent Runs, Sessions, Jobs, Tool Invocations) include a \texttt{tenant\_id} foreign key with database-level constraints.
    
    \item \textbf{Cache Isolation:} Redis keys are namespaced by tenant ID to prevent cache poisoning across tenants.
\end{enumerate}

%------------------------------------------------------------------------------
\subsection{Reliability Requirements}
\label{subsec:nfr-reliability}

Requirements ensuring system availability and fault tolerance:

\begin{table}[htbp]
\centering
\caption{Reliability Non-Functional Requirements}
\label{tab:nfr-reliability}
\small
\begin{tabularx}{\textwidth}{|l|X|l|}
\hline
\textbf{ID} & \textbf{Requirement Description} & \textbf{Priority} \\
\hline
NFR-RE-01 & The system shall implement circuit breakers for all external service calls (LLM providers, databases) with configurable thresholds. & Must \\
\hline
NFR-RE-02 & The system shall support graceful degradation when LLM providers are unavailable, falling back to deterministic responses or cached results. & Must \\
\hline
NFR-RE-03 & The system shall handle Redis unavailability without complete service failure, degrading non-critical features (caching, rate limiting). & Should \\
\hline
NFR-RE-04 & The system shall implement retry logic with exponential backoff for transient failures. & Must \\
\hline
NFR-RE-05 & The system shall support graceful shutdown, completing in-flight requests before termination. & Must \\
\hline
NFR-RE-06 & The system shall maintain heartbeat signals for long-running operations to detect stale workers. & Should \\
\hline
NFR-RE-07 & The system shall achieve 99.5\% availability during normal operations (excluding planned maintenance). & Should \\
\hline
NFR-RE-08 & The system shall recover from transient database connection failures within 30 seconds. & Must \\
\hline
\end{tabularx}
\end{table}

\subsubsection{Circuit Breaker Implementation}

The platform implements a three-state circuit breaker pattern for LLM provider resilience:

\begin{description}
    \item[CLOSED:] Normal operation; requests flow to the provider. Failures are counted.
    \item[OPEN:] Provider is considered unavailable; requests fail fast without attempting the call. Automatic transition to HALF-OPEN after a configurable timeout.
    \item[HALF-OPEN:] A limited number of probe requests are allowed. Success transitions to CLOSED; failure returns to OPEN.
\end{description}

Circuit breaker state is persisted in Redis for cluster-wide consistency, with the following configurable parameters:
\begin{itemize}
    \item \texttt{failure\_threshold}: Number of failures before opening (default: 5)
    \item \texttt{success\_threshold}: Successes in half-open before closing (default: 3)
    \item \texttt{timeout\_seconds}: Duration of open state before half-open (default: 60)
\end{itemize}

%------------------------------------------------------------------------------
\subsection{Performance Requirements}
\label{subsec:nfr-performance}

Latency and throughput requirements for the platform:

\begin{table}[htbp]
\centering
\caption{Performance Non-Functional Requirements}
\label{tab:nfr-performance}
\small
\begin{tabularx}{\textwidth}{|l|X|l|}
\hline
\textbf{ID} & \textbf{Requirement Description} & \textbf{Priority} \\
\hline
NFR-PF-01 & Health probe endpoints (\texttt{/health/live}) shall respond within 10ms (p99). & Must \\
\hline
NFR-PF-02 & API endpoints (excluding LLM operations) shall respond within 500ms (p95). & Should \\
\hline
NFR-PF-03 & Tool invocations shall complete within their configured timeout (default: 30 seconds). & Must \\
\hline
NFR-PF-04 & The system shall support at least 100 concurrent API connections per instance. & Should \\
\hline
NFR-PF-05 & Database queries shall complete within 100ms (p95) for indexed operations. & Should \\
\hline
NFR-PF-06 & Cache hit latency (Redis) shall be under 5ms (p95). & Should \\
\hline
NFR-PF-07 & The system shall implement request-level timeouts to prevent resource exhaustion. & Must \\
\hline
NFR-PF-08 & Background job throughput shall support at least 10 jobs per second per worker. & Should \\
\hline
\end{tabularx}
\end{table}

%------------------------------------------------------------------------------
\subsection{Auditability and Compliance Requirements}
\label{subsec:nfr-audit}

Requirements for audit logging and regulatory compliance:

\begin{table}[htbp]
\centering
\caption{Auditability Non-Functional Requirements}
\label{tab:nfr-audit}
\small
\begin{tabularx}{\textwidth}{|l|X|l|}
\hline
\textbf{ID} & \textbf{Requirement Description} & \textbf{Priority} \\
\hline
NFR-AU-01 & The system shall maintain append-only audit logs for all security-relevant events (authentication, authorization, data access). & Must \\
\hline
NFR-AU-02 & Audit logs shall include: timestamp, principal, action, resource, tenant, outcome, and correlation ID. & Must \\
\hline
NFR-AU-03 & Audit logs shall be retained for a configurable period (default: 90 days) with archival support. & Must \\
\hline
NFR-AU-04 & The system shall support audit log export for compliance reporting and external SIEM integration. & Should \\
\hline
NFR-AU-05 & All agent runs shall record complete execution traces including prompts, responses, and tool calls. & Must \\
\hline
NFR-AU-06 & The system shall track model usage with token counts and estimated costs per tenant. & Should \\
\hline
NFR-AU-07 & Administrative actions (tenant creation, model configuration, permission changes) shall generate audit events. & Must \\
\hline
NFR-AU-08 & The system shall support data residency requirements through configurable deployment regions. & Could \\
\hline
\end{tabularx}
\end{table}

%------------------------------------------------------------------------------
\subsection{Observability Requirements}
\label{subsec:nfr-observability}

Requirements for monitoring, logging, and tracing:

\begin{table}[htbp]
\centering
\caption{Observability Non-Functional Requirements}
\label{tab:nfr-observability}
\small
\begin{tabularx}{\textwidth}{|l|X|l|}
\hline
\textbf{ID} & \textbf{Requirement Description} & \textbf{Priority} \\
\hline
NFR-OB-01 & The system shall expose Prometheus metrics via \texttt{/metrics} endpoint covering HTTP requests, background jobs, tools, and LLM operations. & Must \\
\hline
NFR-OB-02 & The system shall implement structured logging (JSON format) with consistent field schemas. & Must \\
\hline
NFR-OB-03 & The system shall propagate distributed trace context (W3C Trace Context) across all service calls. & Should \\
\hline
NFR-OB-04 & The system shall correlate logs, metrics, and traces using a common request ID. & Must \\
\hline
NFR-OB-05 & Health probes shall check all external dependencies with individual status reporting. & Must \\
\hline
NFR-OB-06 & The system shall provide pre-configured Grafana dashboards for operational monitoring. & Should \\
\hline
NFR-OB-07 & Log output shall be configurable (JSON for production, human-readable for development). & Should \\
\hline
NFR-OB-08 & Metrics shall include latency histograms with configurable bucket boundaries. & Should \\
\hline
\end{tabularx}
\end{table}

\subsubsection{Three Pillars of Observability}

The platform implements comprehensive observability through:

\begin{enumerate}
    \item \textbf{Metrics (Prometheus):} Counter, histogram, and gauge metrics for all major subsystems:
    \begin{itemize}
        \item HTTP request rate, latency, and error rate by endpoint
        \item Background job execution counts and durations
        \item Tool invocation statistics by tool name and outcome
        \item LLM call latency, token consumption, and cost estimates
        \item Rate limit hit rates and quota utilization
    \end{itemize}
    
    \item \textbf{Logging (Structlog):} Structured JSON logs with:
    \begin{itemize}
        \item Consistent field naming across all components
        \item Request ID correlation for distributed tracing
        \item Log level filtering with environment-specific defaults
        \item PII scrubbing before log emission
    \end{itemize}
    
    \item \textbf{Tracing (OpenTelemetry):} Distributed tracing with:
    \begin{itemize}
        \item Automatic span creation for HTTP requests and database calls
        \item Manual instrumentation for LLM calls and tool invocations
        \item Configurable sampling rates (20\% in production, 100\% in development)
        \item OTLP export to compatible collectors
    \end{itemize}
\end{enumerate}

%------------------------------------------------------------------------------
\subsection{Scalability Requirements}
\label{subsec:nfr-scalability}

Requirements for horizontal and vertical scaling:

\begin{table}[htbp]
\centering
\caption{Scalability Non-Functional Requirements}
\label{tab:nfr-scalability}
\small
\begin{tabularx}{\textwidth}{|l|X|l|}
\hline
\textbf{ID} & \textbf{Requirement Description} & \textbf{Priority} \\
\hline
NFR-SC-01 & The system shall support horizontal scaling of API servers behind a load balancer. & Must \\
\hline
NFR-SC-02 & The system shall support horizontal scaling of background workers with job queue partitioning. & Should \\
\hline
NFR-SC-03 & The system shall maintain stateless API servers, storing all session state in external stores (Redis, PostgreSQL). & Must \\
\hline
NFR-SC-04 & The system shall support connection pooling for database connections with configurable pool sizes. & Must \\
\hline
NFR-SC-05 & The system shall implement request-level rate limiting to prevent resource exhaustion during traffic spikes. & Must \\
\hline
NFR-SC-06 & The system architecture shall not preclude deployment on Kubernetes with pod autoscaling. & Should \\
\hline
\end{tabularx}
\end{table}

%------------------------------------------------------------------------------
\subsection{Rate Limiting Requirements}
\label{subsec:nfr-ratelimit}

Requirements for request throttling and quota management:

\begin{table}[htbp]
\centering
\caption{Rate Limiting Non-Functional Requirements}
\label{tab:nfr-ratelimit}
\small
\begin{tabularx}{\textwidth}{|l|X|l|}
\hline
\textbf{ID} & \textbf{Requirement Description} & \textbf{Priority} \\
\hline
NFR-RL-01 & The system shall implement sliding window rate limiting with Redis-backed counters. & Must \\
\hline
NFR-RL-02 & Rate limits shall be configurable per endpoint, per user, and per tenant. & Must \\
\hline
NFR-RL-03 & Rate limit responses shall include \texttt{Retry-After} headers with backoff guidance. & Must \\
\hline
NFR-RL-04 & The system shall support rate limit overrides for privileged users or emergency access. & Should \\
\hline
NFR-RL-05 & Rate limit metrics shall be exposed for monitoring and alerting. & Should \\
\hline
NFR-RL-06 & Default rate limits shall be: 60 requests/minute for general API, 10 requests/minute for LLM operations. & Should \\
\hline
\end{tabularx}
\end{table}

%------------------------------------------------------------------------------
\subsection{Requirements Summary}
\label{subsec:nfr-summary}

Table~\ref{tab:nfr-summary} summarizes the non-functional requirements by category:

\begin{table}[htbp]
\centering
\caption{Non-Functional Requirements Summary by Category}
\label{tab:nfr-summary}
\begin{tabular}{|l|c|c|c|c|}
\hline
\textbf{Category} & \textbf{Must} & \textbf{Should} & \textbf{Could} & \textbf{Total} \\
\hline
Security & 10 & 0 & 0 & 10 \\
Multi-Tenancy & 6 & 2 & 0 & 8 \\
Reliability & 5 & 3 & 0 & 8 \\
Performance & 3 & 5 & 0 & 8 \\
Auditability & 5 & 2 & 1 & 8 \\
Observability & 4 & 4 & 0 & 8 \\
Scalability & 4 & 2 & 0 & 6 \\
Rate Limiting & 3 & 3 & 0 & 6 \\
\hline
\textbf{Total} & \textbf{40} & \textbf{21} & \textbf{1} & \textbf{62} \\
\hline
\end{tabular}
\end{table}

The predominance of \textbf{Must} requirements in the Security category reflects the enterprise context and the sensitive nature of AI agent operations in research environments.
